%! Unit 8: Learning from Data — Summative Assessment — Practice Test --- Study Copy (blank)
\documentclass[10pt]{article}
\usepackage{linalg-article}
\usepackage{linalg-boxes}
\newcommand{\vv}{\mathbf{v}}
\newcommand{\bb}{\mathbf{b}}
\newcommand{\ww}{\mathbf{w}}
\newcommand{\relu}{\mathrm{ReLU}}

% Part-divider strip
\newcommand{\parthead}[1]{%
  \vspace{6pt}\noindent
  \begin{headlinebox}{burgundy}{\color{white}\bfseries #1}\end{headlinebox}%
  \vspace{4pt}}

\begin{document}

\pageheader{Unit 8: Learning from Data}{Practice Test --- Study Copy}
\namedateperiod

\vspace{0.10in}
\begin{remindbox}
\small
\textbf{This is a practice test.} It mirrors the real Unit 8 test in length, format,
and the ideas it covers, but the numbers and contexts differ. Work every problem
with no notes, then check your answers against the key.
\end{remindbox}

% ======================================================================
\parthead{Part A --- Vocabulary (8 pts)}
Write the letter of the best description next to each term.

\vspace{4pt}
\noindent
\begin{minipage}[t]{0.40\textwidth}
\begin{enumerate}[leftmargin=2.2em, itemsep=7pt, label=\arabic*.]
  \item ReLU \blank{1.2cm}
  \item Neural network layer \blank{1.2cm}
  \item Piecewise linear function \blank{1.2cm}
  \item Convolution filter \blank{1.2cm}
  \item Weight sharing \blank{1.2cm}
  \item Gradient descent \blank{1.2cm}
  \item Learning rate \blank{1.2cm}
  \item Covariance matrix \blank{1.2cm}
\end{enumerate}
\end{minipage}%
\hfill
\begin{minipage}[t]{0.57\textwidth}
\small
\begin{description}[leftmargin=1.4em, itemsep=3pt]
  \item[(A)] The building block $\relu(A\vv+\bb)$: multiply by weights, add a bias, then rectify.
  \item[(B)] The step size $s$ in gradient descent --- too small crawls, too big overshoots.
  \item[(C)] The rectifier $\max(0,x)$: keep positive inputs, send negatives to $0$.
  \item[(D)] A small weight pattern slid across the input to detect one feature everywhere.
  \item[(E)] The rule $\ww\leftarrow\ww-s\,\nabla L$: step opposite the slope to shrink the loss.
  \item[(F)] A function built of straight pieces joined at folds (a sum of shifted ReLUs).
  \item[(G)] Reusing one filter at every position, so a giant layer needs only a few weights.
  \item[(H)] The symmetric matrix of variances and covariances; its eigenvectors are the principal components.
\end{description}
\end{minipage}

% ======================================================================
\parthead{Part B --- Multiple Choice (12 pts)}
Circle the best answer.

\begin{enumerate}[leftmargin=1.9em, itemsep=9pt]
  \item $\relu(-3)$ equals
    \hfill (A) $-3$ \quad (B) $0$ \quad (C) $3$ \quad (D) $1$
  \item One neural-network layer computes
    \hfill (A) $A\vv$ \quad (B) $\relu(A\vv+\bb)$ \quad (C) $\vv+\bb$ \quad (D) $A^{-1}\vv$
  \item A piecewise linear function built from $N$ folds (kinks) has
    \hfill (A) $N-1$ pieces \quad (B) $N$ pieces \quad (C) $N+1$ pieces \quad (D) $2N$ pieces
  \item Sliding one small filter across the whole input is convolution; reusing that filter everywhere is
    \hfill (A) weight sharing \quad (B) the bias \quad (C) the loss \quad (D) a rotation
  \item Gradient descent updates a weight by stepping
    \hfill (A) along the slope \quad (B) opposite the slope (downhill) \quad (C) to a random weight \quad (D) to $0$
  \item A positive covariance $\sigma_{xy}>0$ between two features means they
    \hfill (A) move in opposite directions \quad (B) tend to move together \quad (C) are unrelated \quad (D) are equal
\end{enumerate}

% ======================================================================
\parthead{Part C --- Short Answer \& Computation (35 pts)}
Show your work.

\begin{enumerate}[leftmargin=1.9em, itemsep=12pt]
  \item Evaluate one layer $\relu(A\vv+\bb)$ where
        $A=\begin{bsmallmatrix}1&1\\1&-1\end{bsmallmatrix}$, $\vv=\begin{bsmallmatrix}1\\3\end{bsmallmatrix}$,
        $\bb=\begin{bsmallmatrix}0\\1\end{bsmallmatrix}$. Show $A\vv$, then $A\vv+\bb$, then the ReLU.
        \vspace{2.6cm}
  \item The tent function $F(x)=\relu(x)-2\relu(x-2)+\relu(x-4)$.
        \textbf{(a)} Find $F(0),F(2),F(4),F(6)$. \textbf{(b)} Where are its folds, and what is the peak point?
        \vspace{2.8cm}
  \item Slide the edge filter $\ww=[-1,1]$ (output $=$ next $-$ current) across $\vv=[2,2,2,6,6]$.
        \textbf{(a)} Write the feature map. \textbf{(b)} Where is the edge, and why is the output $0$ on the flat parts?
        \vspace{2.6cm}
  \item The filter $[-1,1]$ on a length-$4$ signal is the convolution matrix
        $A=\begin{bsmallmatrix}-1&1&0&0\\0&-1&1&0\\0&0&-1&1\end{bsmallmatrix}$.
        \textbf{(a)} Compute $A\begin{bsmallmatrix}1\\1\\6\\6\end{bsmallmatrix}$.
        \textbf{(b)} How many weights does this shared filter use, versus a full $3\times4$ layer?
        \vspace{2.8cm}
  \item Gradient descent on the loss bowl $L(w)=(w-4)^2$, with slope $L'(w)=2w-8$, learning rate $s=0.25$,
        starting at $w_0=0$. Use $w\leftarrow w-s\,L'(w)$.
        \textbf{(a)} Find $w_1,w_2,w_3$. \textbf{(b)} Give the loss at each step. \textbf{(c)} What weight is it heading toward, and how would you know to stop?
        \vspace{3.0cm}
  \item Two-weight loss $L=(w_1-1)^2+(w_2-3)^2$ has gradient $\nabla L=\big(2(w_1-1),\,2(w_2-3)\big)$.
        Starting at $(0,0)$ with $s=0.25$, take \textbf{one} gradient-descent step. Where does the loss bottom out?
        \vspace{2.8cm}
  \item Four students' two quiz scores are $(2,4),(4,2),(6,8),(8,6)$.
        \textbf{(a)} Find the mean point. \textbf{(b)} Find the variances $\sigma_x^2,\sigma_y^2$ and the covariance $\sigma_{xy}$.
        \textbf{(c)} Write the covariance matrix $V$. \textbf{(d)} Find its eigenvalues and the fraction of the total
        variance the first principal component ($\mathrm{PC}_1$) holds.
        \vspace{3.2cm}
\end{enumerate}

% ======================================================================
\parthead{Part D --- Extended Response (10 pts)}
Explain your reasoning in full sentences.

\begin{enumerate}[leftmargin=1.9em, itemsep=16pt]
  \item A neural network that learns is built entirely from tools in this course.
        \textbf{(a)} Explain what a single layer $\relu(A\vv+\bb)$ does, and why widening it (adding more shifted ReLUs)
        produces a flexible piecewise-linear function that can bend to fit data.
        \textbf{(b)} Explain how gradient descent \emph{chooses} the weights by rolling downhill on the squared-error loss,
        and why a slope of $0$ marks the best weights.
        \vspace{3.0cm}
  \item The covariance matrix $V=\begin{bsmallmatrix}\sigma_x^2&\sigma_{xy}\\\sigma_{xy}&\sigma_y^2\end{bsmallmatrix}$.
        \textbf{(a)} Explain why $V$ is symmetric, and why that guarantees real eigenvalues and perpendicular eigenvectors.
        \textbf{(b)} Explain what those eigenvectors and eigenvalues tell you about the data cloud, and why the total
        variance equals the trace $=$ the sum of the eigenvalues.
        \vspace{2.8cm}
\end{enumerate}

\end{document}
